Este trabajo presenta un sistema de información para el análisis comparativo y predictivo de la evolución de competencias académicas entre las pruebas Saber 11 y Saber Pro. Para el desarrollo del estudio se utilizaron cerca de $1,32$ millones de registros de estudiantes emparejados de manera individual. Con esta información se integraron los resultados de ambas pruebas en una base de datos relacional y se diseñó un modelo de valor agregado apoyado en una red neuronal artificial de tipo perceptrón multicapa. Los resultados obtenidos por este modelo se compararon con los realmente alcanzados mediante técnicas de referencia como la regresión lineal y el algoritmo de agrupamiento K-means. El proceso de estimación se realizó bajo dos escenarios diferentes: uno que incluía variables socioeconómicas y otro que no las consideraba. Esta comparación permitió analizar en qué medida los factores de contexto influyen en el desempeño académico de los estudiantes. Los resultados muestran una capacidad explicativa moderada ($R^2$ entre $0,39$ y $0,61$), en la que la red neuronal supera a los modelos de referencia salvo en inglés, y un valor agregado institucional mayoritariamente negativo. Estos hallazgos no confirman la hipótesis inicial, pero el sistema desarrollado cumple su función analítica y se materializa en un tablero interactivo de apoyo a la toma de decisiones institucionales



\palabrasclave{valor agregado educativo, Saber 11, Saber Pro, redes neuronales artificiales, sistema de información}